\subsection{Zones from the inputs of a task: the first test without instructions}
\label{sec:noinstruction}

This subsection reports the first test of the rule of Section~\ref{sec:bound}: a version of the zones that needs no instruction, with a guarantee on how much can be forgotten. It was run before the recogniser of Section~\ref{sec:recogniser} existed, and what it found is the reason why that recogniser was built; the reader has met the recogniser first, and here the story goes back to this earlier step.
The rule and its guarantee are Theorem~\ref{thm:bound} (Section~\ref{sec:bound}): a zone is built from the inputs of the task itself, its size is set on inputs of the task that were held out, and then, for any model and any later tasks, each input of the task changes its answer with probability at most a chosen level.
The test below is a first, small check of that statement on benchmark files. It is routes only, on a CPU: no model was trained and no GPU was used, so it says nothing yet about answers.

\paragraph{Set-up (fixed in the plan of the repository before the run).} We use the official benchmark files, the instances only; the definition of a task is never read. Three sequences: the 15 tasks of SuperNI order~1 of Section~\ref{sec:protocol}, the 15 tasks of the Long Sequence benchmark in its
official order~3, and the two together as one sequence of \BoundTasksBoth{} tasks. All the training inputs of a task are used for its zone (the plan said 800; the largest file, of the Long Sequence benchmark, has \BoundTrainMax{}), and the official test inputs, which are never used for a zone, for the measures; so the condition of the theorem, that the test inputs are like the training inputs, is assumed here for the benchmark's own split; the coverage line below checks it.
The frozen feature of an input is the weighted, rescaled average of the T5-base word-vectors of all its tokens (cut at 256 tokens), as in the rule \centre{} (its token weights and rescaling use the dev files of ten of the Long Sequence tasks, frozen, not their training or test data). The score $a_k(x)$ of the primary setting is the distance of the feature to its 10th nearest neighbour among the fit inputs of task $k$; a second score, reported next to it, is the distance to the mean of the fit features. A third score, designed later, is described below. When task $k$ arrives,
\BoundMCal{} of its training inputs are held out for calibration (for the \BoundSmallTasks{} tasks with fewer than 400 training inputs, half of them, a deviation from the plan forced by the data; the smallest file has \BoundSmallMin{}), the level is $\beta = 0.05$ (levels 0.1 and 0.2 are reported), and $q_k$ is the split-conformal quantile of Section~\ref{sec:bound}. An input goes to the oldest zone that holds it; an input outside every zone goes to the task with the smallest score divided by its $q_k$.
Four measures, per task and averaged over the tasks of a sequence: the share of its test inputs inside its own zone (covered); the share inside the zone of an older task when the task arrives (blocked); the share given to the own task after the last task (own task); the share whose route after the last task differs from the route when their task arrived (changed).

\paragraph{Two lines, fixed before the run.} (1) \emph{The claim:} on each sequence the average share of changed routes must be at most $\beta$ plus two standard errors and the average coverage at least $1 - \beta$ minus two standard errors; if not, the assumption that test inputs are like training inputs fails on these files, and the paper has to say so. (2) \emph{The use:} ``works without instructions'' on a sequence if on average at least \BoundLineOwn{}\% of the test inputs go to their own task and at most \BoundLineBlock{}\% are blocked; ``does not'' if fewer than \BoundLineNotOwn{}\% go to their own task or more than \BoundLineNotBlock{}\% are blocked; otherwise partial.

\begin{table}[t]
  \caption{Zones from the inputs of a task, no instruction (routes only; $\beta = 0.05$; means over the tasks of a sequence, in \%). Covered: test inputs inside their own zone (promised at least \BoundCovPromised{}\% on average). Answers changed: test inputs whose route after the last task differs from the route when their task arrived (promised at most \BoundAlpha{}\%).
  Blocked: test inputs inside the zone of an older task when their task arrives. Own task: test inputs given to their own task after the last task. ``sentence enc.'': the frozen all-MiniLM-L6-v2 on the instance text, mean pooling. 10-NN: score is the distance to the 10th nearest fit input; mean: distance to the mean of the fit inputs; ``word-vec.'': the word-vector feature.}
  \label{tab:boundpilot}
  \centering
  \footnotesize
  \setlength{\tabcolsep}{3pt}
  \begin{tabular}{lllrrrr}
    \toprule
    Sequence & Feature & Score & Covered, \% & \shortstack[r]{Routes\\changed, \% $\downarrow$} & \shortstack[r]{Blocked inputs,\\\% $\downarrow$} & \shortstack[r]{Sent to own\\task, \% $\uparrow$} \\
    \midrule
    SuperNI (15) & word-vec. & 10-NN & 96.1 & 0.7 & 86.5 & 12.7 \\
    SuperNI (15) & word-vec. & mean & 95.7 & 0.5 & 86.4 & 12.9 \\
    SuperNI (15) & word-vec. & ratio & 94.3 & 0.0 & 65.5 & 33.7 \\
    SuperNI (15) & sentence enc. & 10-NN & 94.7 & 1.4 & 46.6 & 51.0 \\
    SuperNI (15) & sentence enc. & mean & 95.2 & 0.9 & 59.8 & 37.9 \\
    SuperNI (15) & sentence enc. & ratio & 95.9 & 1.8 & 39.9 & 58.0 \\
    \midrule
    Long (15) & word-vec. & 10-NN & 94.9 & 1.1 & 88.3 & 10.6 \\
    Long (15) & word-vec. & mean & 95.0 & 1.0 & 88.6 & 10.4 \\
    Long (15) & word-vec. & ratio & 94.3 & 1.1 & 81.0 & 17.8 \\
    Long (15) & sentence enc. & 10-NN & 93.7 & 3.0 & 53.2 & 43.5 \\
    Long (15) & sentence enc. & mean & 94.0 & 2.9 & 62.6 & 34.4 \\
    Long (15) & sentence enc. & ratio & 94.4 & 1.4 & 54.3 & 43.8 \\
    \midrule
    Both (30) & word-vec. & 10-NN & 95.5 & 0.5 & 93.2 & 6.3 \\
    Both (30) & word-vec. & mean & 95.4 & 0.7 & 93.2 & 6.2 \\
    Both (30) & word-vec. & ratio & 94.3 & 1.5 & 78.1 & 20.3 \\
    Both (30) & sentence enc. & 10-NN & 94.2 & 2.3 & 63.5 & 34.1 \\
    Both (30) & sentence enc. & mean & 94.6 & 1.9 & 75.2 & 22.8 \\
    Both (30) & sentence enc. & ratio & 95.2 & 1.5 & 62.6 & 35.7 \\
    \bottomrule
\end{tabular}

\end{table}

\paragraph{Result.} The claim holds on all three sequences (Table~\ref{tab:boundpilot}). Coverage is \BoundSuperniWordKnnFiveCov{}\%, \BoundLongWordKnnFiveCov{}\% and \BoundBothWordKnnFiveCov{}\% (the line fixed in advance is \BoundSuperniCovLimit{}\%, \BoundLongCovLimit{}\% and \BoundBothCovLimit{}\%), and the share of changed routes is \BoundSuperniWordKnnFiveChanged{}\%, \BoundLongWordKnnFiveChanged{}\% and \BoundBothWordKnnFiveChanged{}\% (the line is \BoundSuperniChangedLimit{}\%, \BoundLongChangedLimit{}\% and \BoundBothChangedLimit{}\%).
None of the test inputs that were inside their own zone when their task arrived changed its route, as Theorem~\ref{thm:writeonce} says. The distance to the mean gives the same picture, and so do the levels 0.1 and 0.2 (on SuperNI: coverage \BoundSuperniWordKnnTenCov{}\% and \BoundSuperniWordKnnTwentyCov{}\%, changed \BoundSuperniWordKnnTenChanged{}\% and \BoundSuperniWordKnnTwentyChanged{}\%).
The use fails, by the lines fixed in advance: \BoundSuperniWordKnnFiveBlocked{}\%, \BoundLongWordKnnFiveBlocked{}\% and \BoundBothWordKnnFiveBlocked{}\% of the test inputs lie inside the zone of an older task when their task arrives, and only \BoundSuperniWordKnnFiveOwn{}\%, \BoundLongWordKnnFiveOwn{}\% and \BoundBothWordKnnFiveOwn{}\% are given to their own task, far from the \BoundLineOwn{}\% and \BoundLineBlock{}\% asked for and beyond the lines for ``does not''.

\paragraph{A second score, designed after the first had failed.} The first score says how unlike the task's fit inputs an input is, and in this feature the inputs of other tasks are about as close as the task's own outliers. What worked for the fallback in Section~\ref{sec:paraphrase} was to measure in units of the distance to unrelated text, and the same idea is used here:
$a_k(x) = d_k(x) / d_0(x)$, where $d_k$ is the distance to the 10th nearest fit input of task $k$ as before and $d_0$ is the distance to the 10th nearest input of a fixed pool of \BoundPool{} unrelated inputs (instances only, no definitions; a few from each of the real tasks of Section~\ref{sec:many} that are not in the sequences; fixed before the first task; for the Long Sequence tasks the pool contains instances of other task files built on the same datasets, for example Yelp, IMDB and MNLI, which were not removed). Everything else is as before, with the same lines.
It was designed after the numbers of the first score had been seen, on the same sequences, and has no free setting; it is therefore a second attempt and not an independent test.

\paragraph{Result for both scores.} The claim holds in every row of Table~\ref{tab:boundpilot} (word-vector feature): the share of changed routes is \BoundAllChangedLow{}\% to \BoundAllChangedHigh{}\% where \BoundAlpha{}\% is allowed, and the coverage is \BoundAllCovLow{}\% to \BoundAllCovHigh{}\% where at least \BoundCovPromised{}\% is promised on average, within two standard errors. The use fails in every row, by the lines fixed in advance (far more than \BoundLineNotBlock{}\% blocked).
The ratio score helps: on the SuperNI tasks \BoundSuperniWordRatioFiveOwn{}\% of the test inputs reach their own task instead of \BoundSuperniWordKnnFiveOwn{}\%, and \BoundToldApartCount{} of the 15 tasks are told apart almost completely (at least \BoundToldApartMin{}\% of their test inputs reach their own task): \texttt{\BoundToldApartNames{}}, that is the three summarisation tasks, the dialogue task and the reading-comprehension task.
It is not enough: nine tasks have at most 10\% of their test inputs given to themselves (\BoundNotApartNames{}), the short-input tasks and the question tasks among them, because their inputs lie inside older zones. With the frozen sentence encoder as the feature (all-MiniLM-L6-v2 on the instance text, mean pooling; rows ``sentence enc.'' of Table~\ref{tab:boundpilot}) the claim holds again (coverage \BoundMiniCovLow{}\% to \BoundMiniCovHigh{}\%, at most \BoundMiniChangedMax{}\% of the routes changed in any row) and the use still fails by the lines fixed in advance (with the ratio score \BoundSuperniMiniRatioFiveBlocked{}\%, \BoundLongMiniRatioFiveBlocked{}\% and \BoundBothMiniRatioFiveBlocked{}\% blocked on the three sequences), but the sentence encoder tells the tasks apart better than the word-vector average: on the SuperNI tasks \BoundSuperniMiniRatioFiveBlocked{}\% blocked instead of \BoundSuperniWordRatioFiveBlocked{}\% and \BoundSuperniMiniRatioFiveOwn{}\% of the test inputs given to their own task instead of \BoundSuperniWordRatioFiveOwn{}\%. It was the secondary feature of the plan, so it changes no verdict.


\paragraph{Pilot B: trained answers on a small model.} To see the same thing with answers, the stand-alone code (\texttt{certified\_switch/}) was run on \texttt{t5-small} with the five tasks of the paper, with the definition removed from every input (training and test), \BoundBTrain{} training and \BoundBTest{} test inputs per task, \BoundBEpochs{} epochs, one seed, on a CPU; \BoundBCal{} calibration inputs per task, $\beta = 0.05$. Rules compared on the same adapters: all adapters on, and the zones with the plain score or with the ratio score.
Lines fixed before the first run: ``works'' if with the zones at most 10\% of the \BoundBCompared{} answers to earlier tasks change and forgetting is at most 3 points, while with all adapters on more than half change; ``does not'' if more than 20\% change. It is one seed and a small model: a demonstration, not a result about scale.

\begin{table}[t]
  \caption{Pilot B, \texttt{t5-small}, five tasks, no definition in any input, one seed. Scores are Rouge-L averages over the tasks: ``Learned'', right after each task was learned, and ``End'', at the end; forgetting is the average over the first four tasks of the drop from right after learning to the end (the last task cannot be forgotten), so it is not the difference of the two columns; ``Changed'': of the \BoundBCompared{} answers to the test inputs of the first four tasks, those that differ at the end. ``Own task'': test inputs given to their own task after the last task. The last column lists, per task in order, the test inputs that lie inside the zone of an older task when the task arrives.}
  \label{tab:boundpilotb}
  \centering
  \footnotesize
  \setlength{\tabcolsep}{3pt}
  \begin{tabular}{lrrrrrl}
    \toprule
    Rule & \shortstack[r]{Score\\right after\\learning $\uparrow$} & Final score $\uparrow$ & \shortstack[r]{Forgetting\\(points) $\downarrow$} & \shortstack[r]{Answers\\changed $\downarrow$} & \shortstack[r]{Sent to own\\task, \% $\uparrow$} & Blocked (of 50) \\
    \midrule
    All adapters on & 64.4 & 9.6 & 68.4 & 200 of 200 & -- & -- \\
    Zones, plain score & 24.1 & 24.1 & 0.0 & 0 of 200 & 21 & 0, 48, 50, 47, 50 \\
    Zones, ratio score & 36.7 & 33.6 & 3.8 & 8 of 200 & 50 & 0, 0, 23, 48, 45 \\
    \bottomrule
\end{tabular}

\end{table}

\emph{First run (plain score).} By the lines written before it, it met them: no answer changed (\BoundBPlainChanged{} of \BoundBCompared{}), forgetting \BoundBPlainFT{}, against \BoundBAllChanged{} of \BoundBCompared{} changed with all adapters on. It was nevertheless empty: \BoundBPlainFirstTrainBlocked{} of the \BoundBTrain{} training inputs of the second task lie inside the zone of the first, nearly all inputs of the later tasks do too, and only the first task was learned; with these zones the later tasks were never served by their own adapters, so there was nothing to forget.
The lines had been badly chosen. Before the second run was read they were corrected: ``works'' needs, in addition, that the tasks are learned (the average score right after learning with the zones at least 90\% of that with all adapters on) and that on average at least 80\% of the test inputs go to their own task, and the other conditions stayed as before (at most 20 of the \BoundBCompared{} answers changed, more than 100 with all adapters on, forgetting at most 3 points).
\emph{Second run (ratio score).} With the zones, \BoundBRatioChanged{} of \BoundBCompared{} answers to earlier tasks changed (4\%, where the bound allows 5\% on average) and none of the \BoundBRatioInside{} test inputs that lay inside their own zone changed; forgetting is \BoundBRatioFT{} points, above the 3 allowed. But the tasks are not learned: the average score right after learning is \BoundBRatioLearned{} against \BoundBAllLearned{} with all adapters on (\BoundBRatioRatioToAll{}\%); only the first two tasks reached, right after learning, the score they have with all adapters on (\BoundBRatioScoreOne{} and \BoundBRatioScoreTwo{}, the latter at the level of guessing under both rules); the third reached \BoundBRatioScoreThree{} against \BoundBAllScoreThree{}, the fourth and fifth \BoundBRatioScoreFour{} and \BoundBRatioScoreFive{} against \BoundBAllScoreFour{} and \BoundBAllScoreFive{}; and only about \BoundBRatioOwn{}\% of the test inputs reach their own task. By the corrected lines the second run does not work either.
The bound held in both runs (\BoundBPlainChanged{} and \BoundBRatioChanged{} of \BoundBCompared{} answers changed; no answer to an input inside its own zone changed). One more finding: without its definition the small model reached only the level of guessing (\BoundBAllGuessScore{}) on \BoundBAllGuess{} of the five tasks even with all adapters on; whether these two tasks are learned better with their definition was not compared task by task in this run.

\paragraph{What this subsection shows.} \emph{Reached:} a guarantee that needs no instruction, is more than ``inside a zone nothing changes'' (each input of a task changes its answer with probability at most a chosen level, for any model and any later tasks), is proved (Theorem~\ref{thm:bound}), and was met in every run (routes on 30 benchmark tasks; answers of a small trained model).
\emph{Not reached by these zones:} a method that works without instructions. With frozen word-vector features the zones of early tasks hold the inputs of later tasks; this is reported when a task arrives, nothing is silent, but those tasks are then not learned. The bound moves the whole difficulty to one thing: a frozen feature in which the inputs of different tasks are apart without an instruction (a sentence encoder helps but does not suffice, \BoundBothMiniRatioFiveBlocked{}\% blocked at best on all 30 tasks; the hidden states of the frozen model itself were tried next, and they are the feature of the recogniser of Section~\ref{sec:recogniser}). The test inputs here are like the training inputs (the benchmark's own split); a reworded or shifted input has no guarantee from this theorem. No GPU run was made for these zones, because no pilot worked.
The recogniser of Section~\ref{sec:recogniser}, reported earlier in this section, was built because of these two findings, the blocked tasks and the need for a better feature. It needs no zone and blocks no task.
